% GENERATED FILE. Edit canonical lessons, then run npm run generate:books.
% canonical-source-hash: fnv1a64:a92ace3f719e79fc
% canonical-lessons: ML-C406-viyojikkuka, ML-C406-tarkkam, ML-C406-samvadam, ML-C406-bodhyappetuttuka, ML-C406-viksanam

\chapter{Agreeing and Disagreeing}
\label{ch:ml-agreeing-and-disagreeing}
\hlchaptermodality{406}

\begin{chapteropening}
\textbf{By the end of this chapter:} \emph{I can say to disagree, a dispute, a debate, to convince and an outlook.}

Complete the last lesson of chapter 406: I can say to disagree, a dispute, a debate, to convince and an outlook.
\end{chapteropening}

\section[\texorpdfstring{\ml{വിയോജിക്കുക} (viyōjikkuka)}{വിയോജിക്കുക (viyōjikkuka)}]{\ml{വിയോജിക്കുക} (viyōjikkuka) — to disagree}

\label{lesson:ML-C406-viyojikkuka}



\begin{quote}
Before the new one: say the Malayalam for to admit, then the Malayalam for a conclusion.
\end{quote}

\subsection*{You'll want to know: \ml{വിയോജിക്കുക}}
\textbf{\ml{വിയോജിക്കുക}} — \emph{viyōjikkuka} — \textquotedblleft{}to disagree\textquotedblright{}.

\begin{quote}
\ml{ഞാൻ} \ml{നിങ്ങളോട്} \ml{വിയോജിക്കുന്നു}.
\end{quote}

\begin{quote}
\emph{ñān niṅṅaḷōṭŭ viyōjikkunnu.}
\end{quote}

\begin{quote}
I disagree with you.
\end{quote}

Past: \ml{വിയോജിച്ചു}, \emph{viyōjiccu}.

The person takes -\ml{ഓട്}: \ml{അവനോട്} \ml{വിയോജിക്കുക}. Its opposite is \ml{യോജിക്കുക}, to agree.

\begin{cousinweb}
\textbf{\ml{വിയോജിക്കുക}} (\emph{viyōjikkuka}) is built from Sanskrit \emph{viyojana}, \textquotedblleft{}detaching, separation\textquotedblright{} (Monier-Williams): the prefix \emph{vi-}, \textquotedblleft{}apart\textquotedblright{}, on the root \emph{yuj}, \textquotedblleft{}to join\textquotedblright{}, the root of \ml{യോജിക്കുക}. To disagree is to come apart.
\end{cousinweb}

\begin{grammarlens}[title={What you've built}]
One more word for giving reasons and opinions.
\end{grammarlens}

\subsection*{Your turn}
\begin{itemize}
\raggedright
  \item \emph{Say it:} \emph{viyōjikkuka}
  \item \emph{Say it:} \emph{viyōjikkuka}, once more
  \item \emph{Say it:} \emph{ñān niṅṅaḷōṭŭ viyōjikkunnu.} — I disagree with you.
  \item \emph{Recall:} say the Malayalam for to admit, then the Malayalam for a conclusion, then say \emph{viyōjikkuka} again
\end{itemize}

\begin{tcolorbox}[breakable,colback=teal!4,colframe=teal!35!black,title={Before you move on}]
What does \ml{വിയോജിക്കുക} mean? (\textquotedblleft{}To disagree\textquotedblright{}.) Say it once more.
\end{tcolorbox}
\section[\texorpdfstring{\ml{തർക്കം} (tarkkaṁ)}{തർക്കം (tarkkaṁ)}]{\ml{തർക്കം} (tarkkaṁ) — a dispute; a quarrel}

\label{lesson:ML-C406-tarkkam}



\begin{quote}
Before the new one: say the Malayalam for a conclusion, then the Malayalam for to disagree.
\end{quote}

\subsection*{You'll want to know: \ml{തർക്കം}}
\textbf{\ml{തർക്കം}} — \emph{tarkkaṁ} — \textquotedblleft{}a dispute; a quarrel\textquotedblright{}.

\begin{quote}
\ml{വീടിന്റെ} \ml{വിലയെക്കുറിച്ച്} \ml{അവർ} \ml{തമ്മിൽ} \ml{തർക്കം} \ml{ഉണ്ടായി}.
\end{quote}

\begin{quote}
\emph{vīṭinṟe vilayekkuṟiccŭ avar tammil tarkkaṁ uṇṭāyi.}
\end{quote}

\begin{quote}
There was a dispute between them about the price of the house.
\end{quote}

The verb \ml{തർക്കിക്കുക}, to argue, is the same word.

\begin{cousinweb}
\textbf{\ml{തർക്കം}} (\emph{tarkkaṁ}) is borrowed from Sanskrit \emph{tarka}, \textquotedblleft{}inquiry, conjecture, reasoning\textquotedblright{} (CDIAL 5714). In Malayalam the reasoning became the quarrel it can lead to.
\end{cousinweb}

\begin{grammarlens}[title={What you've built}]
One more word for giving reasons and opinions.
\end{grammarlens}

\subsection*{Your turn}
\begin{itemize}
\raggedright
  \item \emph{Say it:} \emph{tarkkaṁ}
  \item \emph{Say it:} \emph{tarkkaṁ}, once more
  \item \emph{Say it:} \emph{vīṭinṟe vilayekkuṟiccŭ avar tammil tarkkaṁ uṇṭāyi.} — There was a dispute between them about the price of the house.
  \item \emph{Recall:} say the Malayalam for a conclusion, then the Malayalam for to disagree, then say \emph{tarkkaṁ} again
\end{itemize}

\begin{tcolorbox}[breakable,colback=teal!4,colframe=teal!35!black,title={Before you move on}]
What does \ml{തർക്കം} mean? (\textquotedblleft{}A dispute; a quarrel\textquotedblright{}.) Say it once more.
\end{tcolorbox}
\section[\texorpdfstring{\ml{സംവാദം} (saṁvādaṁ)}{സംവാദം (saṁvādaṁ)}]{\ml{സംവാദം} (saṁvādaṁ) — a debate}

\label{lesson:ML-C406-samvadam}



\begin{quote}
Before the new one: say the Malayalam for to disagree, then the Malayalam for a dispute.
\end{quote}

\subsection*{You'll want to know: \ml{സംവാദം}}
\textbf{\ml{സംവാദം}} — \emph{saṁvādaṁ} — \textquotedblleft{}a debate\textquotedblright{}.

\begin{quote}
\ml{കോളേജിൽ} \ml{ഇന്നലെ} \ml{ഒരു} \ml{സംവാദം} \ml{നടന്നു}.
\end{quote}

\begin{quote}
\emph{kōḷējil innale oru saṁvādaṁ naṭannu.}
\end{quote}

\begin{quote}
A debate was held at the college yesterday.
\end{quote}

A debate is \ml{നടക്കുക}, held, or \ml{നടത്തുക}, organised.

\begin{cousinweb}
\textbf{\ml{സംവാദം}} (\emph{saṁvādaṁ}) is borrowed from Sanskrit \emph{saṁvāda}, \textquotedblleft{}speaking together, a conversation\textquotedblright{} (Monier-Williams): \emph{sam-}, \textquotedblleft{}together\textquotedblright{}, and \emph{vāda}, \textquotedblleft{}speech\textquotedblright{}.
\end{cousinweb}

\begin{grammarlens}[title={What you've built}]
One more word for giving reasons and opinions.
\end{grammarlens}

\subsection*{Your turn}
\begin{itemize}
\raggedright
  \item \emph{Say it:} \emph{saṁvādaṁ}
  \item \emph{Say it:} \emph{saṁvādaṁ}, once more
  \item \emph{Say it:} \emph{kōḷējil innale oru saṁvādaṁ naṭannu.} — A debate was held at the college yesterday.
  \item \emph{Recall:} say the Malayalam for to disagree, then the Malayalam for a dispute, then say \emph{saṁvādaṁ} again
\end{itemize}

\begin{tcolorbox}[breakable,colback=teal!4,colframe=teal!35!black,title={Before you move on}]
What does \ml{സംവാദം} mean? (\textquotedblleft{}A debate\textquotedblright{}.) Say it once more.
\end{tcolorbox}
\section[\texorpdfstring{\ml{ബോധ്യപ്പെടുത്തുക} (bōdhyappeṭuttuka)}{ബോധ്യപ്പെടുത്തുക (bōdhyappeṭuttuka)}]{\ml{ബോധ്യപ്പെടുത്തുക} (bōdhyappeṭuttuka) — to convince}

\label{lesson:ML-C406-bodhyappetuttuka}



\begin{quote}
Before the new one: say the Malayalam for a dispute, then the Malayalam for a debate.
\end{quote}

\subsection*{You'll want to know: \ml{ബോധ്യപ്പെടുത്തുക}}
\textbf{\ml{ബോധ്യപ്പെടുത്തുക}} — \emph{bōdhyappeṭuttuka} — \textquotedblleft{}to convince\textquotedblright{}.

\begin{quote}
\ml{ഒടുവിൽ} \ml{അവൾ} \ml{അച്ഛനെ} \ml{കാര്യം} \ml{ബോധ്യപ്പെടുത്തി}.
\end{quote}

\begin{quote}
\emph{oṭuvil avaḷ acchane kāryaṁ bōdhyappeṭutti.}
\end{quote}

\begin{quote}
At last she convinced her father of the matter.
\end{quote}

Past: \ml{ബോധ്യപ്പെടുത്തി}, \emph{bōdhyappeṭutti}.

The person convinced takes -\ml{നെ}, the matter stays unmarked: \ml{അച്ഛനെ} \ml{കാര്യം} \ml{ബോധ്യപ്പെടുത്തുക}.

\begin{cousinweb}
\textbf{\ml{ബോധ്യപ്പെടുത്തുക}} (\emph{bōdhyappeṭuttuka}) is built from Sanskrit \emph{bodhya}, \textquotedblleft{}to be known or understood\textquotedblright{} (Monier-Williams), with -\ml{പ്പെടുത്തുക}, \textquotedblleft{}to cause to be\textquotedblright{}, from native \ml{പെടുക} (DEDR 3853): to make a thing understood.
\end{cousinweb}

\begin{grammarlens}[title={What you've built}]
One more word for giving reasons and opinions.
\end{grammarlens}

\subsection*{Your turn}
\begin{itemize}
\raggedright
  \item \emph{Say it:} \emph{bōdhyappeṭuttuka}
  \item \emph{Say it:} \emph{bōdhyappeṭuttuka}, once more
  \item \emph{Say it:} \emph{oṭuvil avaḷ acchane kāryaṁ bōdhyappeṭutti.} — At last she convinced her father of the matter.
  \item \emph{Recall:} say the Malayalam for a dispute, then the Malayalam for a debate, then say \emph{bōdhyappeṭuttuka} again
\end{itemize}

\begin{tcolorbox}[breakable,colback=teal!4,colframe=teal!35!black,title={Before you move on}]
What does \ml{ബോധ്യപ്പെടുത്തുക} mean? (\textquotedblleft{}To convince\textquotedblright{}.) Say it once more.
\end{tcolorbox}
\section[\texorpdfstring{\ml{വീക്ഷണം} (vīkṣaṇaṁ)}{വീക്ഷണം (vīkṣaṇaṁ)}]{\ml{വീക്ഷണം} (vīkṣaṇaṁ) — an outlook; a perspective}

\label{lesson:ML-C406-viksanam}



\begin{quote}
Before the new one: say the Malayalam for a debate, then the Malayalam for to convince.
\end{quote}

\subsection*{You'll want to know: \ml{വീക്ഷണം}}
\textbf{\ml{വീക്ഷണം}} — \emph{vīkṣaṇaṁ} — \textquotedblleft{}an outlook; a perspective\textquotedblright{}.

\begin{quote}
\ml{ഓരോരുത്തർക്കും} \ml{ഓരോ} \ml{വീക്ഷണം} \ml{ഉണ്ട്}.
\end{quote}

\begin{quote}
\emph{ōrōruttarkkuṁ ōrō vīkṣaṇaṁ uṇṭŭ.}
\end{quote}

\begin{quote}
Each person has their own outlook.
\end{quote}

From the point of view of is -\ml{ന്റെ} \ml{വീക്ഷണത്തിൽ}.

\begin{cousinweb}
\textbf{\ml{വീക്ഷണം}} (\emph{vīkṣaṇaṁ}) is borrowed from Sanskrit \emph{vīkṣaṇa}, \textquotedblleft{}looking at, seeing, inspection\textquotedblright{} (Monier-Williams). An outlook is the way a person looks at things.
\end{cousinweb}

\begin{grammarlens}[title={What you've built}]
One more word for giving reasons and opinions.
\end{grammarlens}

\subsection*{Your turn}
\begin{itemize}
\raggedright
  \item \emph{Say it:} \emph{vīkṣaṇaṁ}
  \item \emph{Say it:} \emph{vīkṣaṇaṁ}, once more
  \item \emph{Say it:} \emph{ōrōruttarkkuṁ ōrō vīkṣaṇaṁ uṇṭŭ.} — Each person has their own outlook.
  \item \emph{Recall:} say the Malayalam for a debate, then the Malayalam for to convince, then say \emph{vīkṣaṇaṁ} again
\end{itemize}

\begin{tcolorbox}[breakable,colback=teal!4,colframe=teal!35!black,title={Before you move on}]
What does \ml{വീക്ഷണം} mean? (\textquotedblleft{}An outlook; a perspective\textquotedblright{}.) Say it once more.
\end{tcolorbox}
